\documentclass[11pt]{article}
\usepackage[margin=1in]{geometry}
\usepackage{times}
\usepackage[T1]{fontenc}
\usepackage{microtype}
\usepackage{enumitem}
\usepackage{titlesec}
\usepackage[hidelinks]{hyperref}
\usepackage{parskip}

\titleformat{\section}{\normalfont\large\bfseries}{\thesection}{0.6em}{}
\titleformat{\subsection}{\normalfont\normalsize\bfseries}{\thesubsection}{0.6em}{}
\setlist[itemize]{leftmargin=1.4em,itemsep=3pt,topsep=3pt}

\title{\vspace{-2em}\textbf{The State of AI-Readiness on Business Websites (2026)}\\[0.4em]
\large A large-scale measurement of AI-crawler permission, llms.txt adoption,\\
and structured-data readiness}

\author{Joshua R. Gutierrez\\
\normalsize Axion Deep Labs, Inc., Las Cruces, NM\\
\normalsize \texttt{hello@axiondeep.com}}

\date{Working paper. Version: July 4, 2026.}

\begin{document}
\maketitle

\begin{abstract}
\noindent As AI answer engines become a discovery channel, a website's ability to participate is
gated by three public, machine-checkable signals: whether it permits AI crawlers (robots.txt),
whether it guides them (llms.txt), and whether it is structured for machine identification
(schema.org). We define AI-readiness operationally as adoption of these three public signals, and
measure them across a stratified sample of local-business websites ($n=766$, 10 verticals, 15
cities) drawn from OpenStreetMap. Among reachable sites, adoption of the structured-data signals
designed to make a business machine-identifiable is low: review/AggregateRating markup appears on
9\% of homepages (95\% CI 7--11) and FAQ markup on 4\% (3--6). The modal business scores 2 of 10 on
a pre-registered AI-Readiness index (42\% of reachable sites): it does not block AI crawlers, but
exposes no schema, reviews markup, or llms.txt. Readiness varies sharply by vertical: accounting
firms show 0\% review and 0\% Person markup (0 of 27; Wilson 95\% CI 0--12 for each),
non-overlapping with legal firms (review 26\%, CI 18--35; Person 31\%, CI 22--41); accounting is the
only vertical simultaneously at zero review, zero Person, and lowest identity markup. llms.txt
adoption is 25\% (22--29), of which at least 28\% carries a tool-generation signature rather than
being authored. Unlike publishers, local businesses rarely block AI crawlers (4\%, CI 3--6). Figures
replicated across 149-, 487-, and 766-site samples. The dataset, code, and codebooks are released
openly.
\end{abstract}

\vspace{0.5em}
\noindent\textbf{Keywords:} AI search, answer engines, generative engine optimization, llms.txt,
schema.org, structured data, AI crawlers, robots.txt, local business, measurement study

\noindent\textbf{JEL classification:} L86, M31, O33

\section{Introduction}

Search interfaces increasingly include AI-generated summaries and conversational answer systems
(ChatGPT, Perplexity, Gemini, Google AI Overview) alongside traditional ranked results. Appearing in
those answers plausibly depends on signals different from classic ranking: whether the engine is
permitted to use the content, whether it can reach it, and whether the business is described in a
machine-readable form. This paper measures the supply side of that question: what businesses
actually expose, before any question is ever asked. It does not measure whether or how AI engines
use these signals; a companion study measures the demand side (what AI engines recommend and cite).

\subsection{Scope of ``AI-readiness''}

We use the term throughout as an \emph{operational construct}, not a comprehensive audit. A full
account of AI-readiness would also include crawlability, rendering, content quality, authority,
citations, reviews on third-party platforms, Google Business Profile completeness, and performance.
We deliberately restrict measurement to three signals that are public, machine-checkable at scale,
and cost nothing to observe: AI-crawler permission (robots.txt), llms.txt, and homepage schema.org
markup. Where we write ``AI-readiness,'' we mean adoption of these three signals, and nothing broader
is claimed.

\medskip
\noindent\textbf{Contributions.} (1) The first stratified, multi-vertical measurement of AI-crawler
permission, llms.txt adoption, and schema readiness together; (2) a pre-registered operational
AI-Readiness index; (3) an open, reproducible, zero-cost methodology and dataset.

\section{Background and related work}

\begin{itemize}
\item \textbf{AI crawlers and permission.} Major model and answer engines publish named crawler
user-agents (GPTBot, Google-Extended, ClaudeBot, PerplexityBot, CCBot, and others) that sites may
allow or disallow in robots.txt. Publisher blocking has been widely reported; local-business
behavior has not been measured systematically.
\item \textbf{llms.txt.} A proposed convention for a root-level markdown file guiding LLMs to a
site's key content. Adoption has not been quantified at scale, nor has the extent of plugin
auto-generation.
\item \textbf{Structured data.} schema.org markup lets machines identify entities, people, and
reviews; its role in AI recommendation is an active question. Prior work (the author's AI-visibility
study) measured whether crawlers can render content; this study measures whether sites are structured
and permitted for AI at all.
\item \textbf{Sampling frames.} Tranco provides a manipulation-resistant top-sites list;
OpenStreetMap provides categorized, geographic local-business coverage used here for the local frame.
\end{itemize}

\section{Methods}

Full protocol and pre-registration: \texttt{PROTOCOL.md}, \texttt{osf\_prereg.md}.

\begin{itemize}
\item \textbf{Frame.} OpenStreetMap POIs carrying a \texttt{website} tag, selected by vertical
(local-service categories) and five geographic strata (metro, mid, small, tourism, university).
\item \textbf{Measurement (public HTTP only).} Per site: robots.txt (AI-crawler permission, versioned
bot list); /llms.txt (validated presence: not an HTML/redirect stub, $\geq 40$ bytes, markdown
structure; plus provenance signature); homepage JSON-LD schema.org \texttt{@type} extraction;
platform fingerprint.
\item \textbf{AI-Readiness index (pre-registered, 0--10):} allows AI crawlers $+2$; valid llms.txt
$+2$; identity schema $+2$; review schema $+2$; FAQ schema $+1$; Person schema $+1$. The index is an
\emph{operational benchmark}, not a validated latent construct: the weights encode a simple editorial
judgment (permission and the three highest-value structured signals count double; the two secondary
schema types count single) and are not derived from factor analysis or an external criterion. It is
intended as a transparent, reproducible summary score for tracking adoption over time, and should not
be read as a psychometric measurement instrument. We report the full component distribution so
readers can reweight.
\item \textbf{Statistics.} Adoption computed over reachable sites; reachability reported separately.
Every proportion (overall and by vertical/stratum) is reported with a Wilson score 95\% confidence
interval; Wilson is used in preference to the percentile bootstrap because several verticals sit at a
boundary proportion (0\% adoption), where the bootstrap degenerates to a zero-width interval and the
Wilson interval remains correctly bounded. The index \emph{mean}, being a mean of 0--10 scores rather
than a proportion, is reported with a bootstrap 95\% CI ($\geq 2000$ resamples, site as unit, fixed
seed). Between-vertical differences are read from non-overlapping 95\% intervals; the flagship
accounting-vs-legal contrast is confirmed by non-overlap on both review and Person markup.
\end{itemize}

\section{Results ($n=766$; reachable $n=594$)}

\noindent\textbf{Reachability.} 78\% of OSM-listed business sites returned a usable homepage (95\% CI
74--80); the remainder were stale URLs, redirects, or blocked. Adoption figures are over reachable
sites.

\medskip
\noindent\textbf{Overall adoption (Wilson 95\% CI).}
\begin{itemize}
\item Identity schema: 49\% (45--53)
\item Review / AggregateRating schema: 9\% (7--11)
\item FAQ schema: 4\% (3--6)
\item Person schema: 10\% (8--13)
\item Valid llms.txt: 25\% (22--29)
\item Blocks at least one AI crawler: 4\% (3--6)
\end{itemize}

\noindent\textbf{The AI-Readiness index.} Mean 3.73 (bootstrap 95\% CI 3.58--3.89), median 4. 42\% of
reachable sites score exactly 2 of 10: they do not block AI crawlers but expose no schema, reviews
markup, or llms.txt.

\medskip
\noindent\textbf{Vertical gradient (primary comparison).} Review-markup prevalence differs sharply by
vertical. Three verticals sit at 0\% review markup (accounting 0/27, optician 0/30, HVAC 0/12), but
accounting is distinctive: it is the only vertical \emph{simultaneously} at 0\% review, 0\% Person,
and the lowest identity markup (26\%, CI 13--45). Its review and Person intervals (0 of 27; Wilson
0--12 for each) do not overlap legal firms (review 26\%, CI 18--35; Person 31\%, CI 22--41), the
flagship contrast. Dentists sit between (review 10\%, CI 6--17).

\medskip
\noindent\emph{Effect size.} Because accounting's review rate is exactly zero, a rate ratio against
it is undefined; we report the contrast as an absolute gap of 26 percentage points with
non-overlapping 95\% intervals. For a stable ratio among sizable verticals, legal firms are
$2.6\times$ more likely than dentists (26\% vs 10\%) and roughly $2.9\times$ the all-vertical review
rate (26\% vs 9\%) to expose review markup. Across ten verticals, the readiness gradient does not
track the value of the services sold; one candidate explanation, developed in the Discussion, is the
maturity of each field's digital-marketing ecosystem, but the design does not adjudicate among
alternatives.

\medskip
\noindent\textbf{llms.txt provenance.} Of the 151 sites with a valid llms.txt, \emph{at least} 28\%
carry an explicit auto-generation signature (Rank Math, Yoast, or other ``generated by'' markers).
This is a lower bound: the remaining 72\% carry no signature but could equally be plugin output
without a marker, copied templates, or agency-supplied files, none of which our detector
distinguishes from hand-authored files. The finding shows that a substantial, and possibly much
larger, share of llms.txt adoption is tool-mediated rather than a deliberate optimization choice; the
exact split between authored and generated files is not identified.

\medskip
\noindent\textbf{Permission.} AI-crawler blocking is rare among local businesses (4\%), in contrast
to the near-universal blocking observed on major news publishers during instrument validation.

\section{Discussion}

\subsection{Permission is not the constraint; structure is}
Within the three signals we measured, the limiting factor is not access but description. Local
businesses overwhelmingly permit AI crawlers (96\% allow all) yet expose little machine-readable
structure: half publish no identity markup and nine in ten publish no review markup. We frame this as
a description gap rather than a recommendation gap on purpose. We did not observe any AI engine's
behavior, so we cannot say these sites are less likely to be recommended; we can say only that they
hand a crawler fewer of the public, machine-readable signals designed to identify and characterize a
business. Whether those signals causally affect AI recommendation is the question our companion
demand-side study takes up, and is out of scope here.

\subsection{Schema is one input among several}
The description gap we measure is specific to on-page structured data. AI systems can also draw on
rendered page text, inbound links and citations, third-party review platforms, knowledge graphs, and
Google Business Profile and Bing Places listings, none of which this study observes. A business with
no homepage schema is therefore not necessarily invisible to AI; it is missing one class of signal it
fully controls and can add at low cost. We caution against reading the schema gap as the sole or even
the primary bottleneck for AI visibility; it is the machine-readable, self-published layer, measured
because it is public and uniform, not because it is established as dominant.

\subsection{Why the vertical gradient, and why accounting}
The gradient does not track the value of the services sold: accounting, among the highest-value
advisory verticals, is the least structured. One plausible explanation is that adoption tracks the
maturity of a vertical's digital-marketing ecosystem: legal marketing has industrialized schema
(specialist agencies, review-driven lead gen, template ecosystems), and accounting has not. But this
is one hypothesis among several the design cannot separate. Accounting firms may on average run older
sites or CMSs, depend more on referral and relationship channels than on search, spend less on SEO,
or skew smaller than the legal firms in our frame; any of these would produce the same pattern
without a ``marketing maturity'' story. Our data support the \emph{observation} of the gap far more
strongly than any single \emph{explanation} for it; disentangling firm age, size, CMS, and channel
mix would require firmographic covariates we did not collect.

\subsection{llms.txt adoption is partly an artifact of defaults}
At least 28\% of llms.txt files are tool-generated, and the true share is plausibly higher. This
reframes headline ``adoption'' as partly a byproduct of plugin defaults rather than deliberate
intent, and is a general caution: raw adoption counts for any emerging convention conflate deliberate
optimization with software that ships the feature by default. Longitudinal tracking that separates
authored from generated files would measure intent more faithfully than a single cross-sectional
count.

\subsection{Implications for AI-visibility practice}
For practitioners, the actionable part is concrete and low-cost: identity, review, and Person markup
plus a maintained llms.txt move a business from the modal 2/10 toward the top of the observed
distribution with no change to permission or content strategy. We frame this as improving
machine-readable self-description, a necessary-looking but not proven-sufficient condition for AI
visibility. Whether it moves actual AI recommendation is exactly what a controlled demand-side study
should test next, and is the natural continuation of this work.

\section{Limitations}

Homepage-only schema (interior pages may differ); single snapshot (adoption is fast-moving; a
longitudinal arm is planned); robots.txt reflects stated policy, not enforcement; platform detection
is heuristic (a share remain unclassified); the OSM frame under-represents rural and tourism
businesses (a real signal, but it thins those strata); provenance detection catches only explicit
signatures. Localization is scoped to US English in this wave. We collected no firmographic
covariates (firm age, size, CMS vintage, marketing spend, channel mix), so the vertical gradient is
reported as an observed association and the design cannot attribute it to any single cause. Most
fundamentally, the study measures the supply of public signals, not their effect: it does not observe
AI-engine behavior, so no claim about actual recommendation or citation is tested here.

\section{Conclusion}

Local businesses are permitted-but-undescribed: they let AI crawlers in and expose little
machine-readable structure. Measured against the three public signals we define as AI-readiness, the
gap is widest in high-value advisory verticals such as accounting. Because the missing signals are
cheap, public, and machine-checkable, the gap is closable, and this study provides a baseline and an
open operational index to track it. Whether closing it changes what AI engines actually recommend is
the demand-side question we take up next.

\section*{Open materials}

All materials are released at \url{https://github.com/Axion-Deep-Labs/ai-readiness-2026}: the query
frame and parsed dataset (149-, 487-, and 766-site crawls; CC BY 4.0), the code (build/crawl/analyze;
MIT), the full protocol and preregistration, and the versioned AI-crawler and schema codebooks. Raw
HTML is not redistributed. All reported figures reproduce with \texttt{python3 analyze.py
data/crawl\_confirmatory\_766.jsonl sample\_meta.csv}.

\begin{thebibliography}{9}
\bibitem{tranco} Le Pochat, V., Van Goethem, T., Tajalizadehkhoob, S., Korczy\'{n}ski, M., \&
Joosen, W. (2019). \emph{Tranco: A Research-Oriented Top Sites Ranking Hardened Against
Manipulation.} Proceedings of the 26th Network and Distributed System Security Symposium (NDSS).
\bibitem{schema} schema.org. \emph{Schema.org vocabulary.} \url{https://schema.org} (accessed 2026).
\bibitem{llmstxt} Howard, J., \& Answer.AI. (2024). \emph{The /llms.txt proposal.}
\url{https://llmstxt.org} (accessed 2026).
\bibitem{osm} OpenStreetMap contributors. \emph{OpenStreetMap.} \url{https://www.openstreetmap.org}
(accessed 2026).
\bibitem{wilson} Wilson, E. B. (1927). \emph{Probable inference, the law of succession, and
statistical inference.} Journal of the American Statistical Association, 22(158), 209--212.
\bibitem{efron} Efron, B., \& Tibshirani, R. J. (1993). \emph{An Introduction to the Bootstrap.}
Chapman \& Hall.
\bibitem{visibility} Gutierrez, J. R. (2026). \emph{Can AI Crawlers See Your Site? A Rendering-Based
AI-Visibility Study.} Axion Deep Labs working paper (companion, supply-side rendering).
\bibitem{recs} Gutierrez, J. R. (2026). \emph{A Large-Scale Analysis of Commercial Recommendations by
AI Search Engines.} Axion Deep Labs working paper (companion, demand-side; in preparation).
\end{thebibliography}

\end{document}
